\documentclass[10pt,a4paper]{amsart}
\usepackage{amsmath,amssymb,amsthm,graphicx}
% Retained minimal preamble for the Stacks Project chapter more-algebra.tex (More on Algebra).
% The chapter file pinned-original-source.tex carries only \input{preamble} on its first line, so no
% preamble text lives inside the chapter itself.  The Stacks Project uses one single shared preamble for
% every chapter; the snapshot stages it identically as sources/stacks-chow/preamble.tex and
% sources/stacks-etale-cohomology/preamble.tex.  The block below is copied mechanically, line for line,
% from sources/stacks-chow/preamble.tex lines 14-18 (the Xy-pic setup) and 154-217 (the native theorem
% declarations together with the mathematical macros).  Nothing was added, removed or redefined; the
% wrapper class and the amsmath/amssymb/amsthm/graphicx base packages are supplied by the assembler.
% Chapter: more-algebra.tex -- More on Algebra.
% For commutative diagrams we use Xy-pic
\usepackage[all]{xy}

% We use 2cell for 2-commutative diagrams.
\xyoption{2cell}

% Theorem environments.
%
\theoremstyle{plain}
\newtheorem{theorem}[subsection]{Theorem}
\newtheorem{proposition}[subsection]{Proposition}
\newtheorem{lemma}[subsection]{Lemma}

\theoremstyle{definition}
\newtheorem{definition}[subsection]{Definition}
\newtheorem{example}[subsection]{Example}
\newtheorem{exercise}[subsection]{Exercise}
\newtheorem{situation}[subsection]{Situation}

\theoremstyle{remark}
\newtheorem{remark}[subsection]{Remark}
\newtheorem{remarks}[subsection]{Remarks}

\numberwithin{equation}{subsection}

% Macros
%
\def\lim{\mathop{\mathrm{lim}}\nolimits}
\def\colim{\mathop{\mathrm{colim}}\nolimits}
\def\Spec{\mathop{\mathrm{Spec}}}
\def\Hom{\mathop{\mathrm{Hom}}\nolimits}
\def\Ext{\mathop{\mathrm{Ext}}\nolimits}
\def\SheafHom{\mathop{\mathcal{H}\!\mathit{om}}\nolimits}
\def\SheafExt{\mathop{\mathcal{E}\!\mathit{xt}}\nolimits}
\def\Sch{\mathit{Sch}}
\def\Mor{\mathop{\mathrm{Mor}}\nolimits}
\def\Ob{\mathop{\mathrm{Ob}}\nolimits}
\def\Sh{\mathop{\mathit{Sh}}\nolimits}
\def\NL{\mathop{N\!L}\nolimits}
\def\CH{\mathop{\mathrm{CH}}\nolimits}
\def\proetale{{pro\text{-}\acute{e}tale}}
\def\etale{{\acute{e}tale}}
\def\QCoh{\mathit{QCoh}}
\def\Ker{\mathop{\mathrm{Ker}}}
\def\Im{\mathop{\mathrm{Im}}}
\def\Coker{\mathop{\mathrm{Coker}}}
\def\Coim{\mathop{\mathrm{Coim}}}

% Boxtimes
%
\DeclareMathSymbol{\boxtimes}{\mathbin}{AMSa}{"02}

%
% Macros for moduli stacks/spaces
%
\def\QCohstack{\mathcal{QC}\!\mathit{oh}}
\def\Cohstack{\mathcal{C}\!\mathit{oh}}
\def\Spacesstack{\mathcal{S}\!\mathit{paces}}
\def\Quotfunctor{\mathrm{Quot}}
\def\Hilbfunctor{\mathrm{Hilb}}
\def\Curvesstack{\mathcal{C}\!\mathit{urves}}
\def\Polarizedstack{\mathcal{P}\!\mathit{olarized}}
\def\Complexesstack{\mathcal{C}\!\mathit{omplexes}}
% \Pic is the operator that assigns to X its picard group, usage \Pic(X)
% \Picardstack_{X/B} denotes the Picard stack of X over B
% \Picardfunctor_{X/B} denotes the Picard functor of X over B
\def\Pic{\mathop{\mathrm{Pic}}\nolimits}
\def\Picardstack{\mathcal{P}\!\mathit{ic}}
\def\Picardfunctor{\mathrm{Pic}}
\def\Deformationcategory{\mathcal{D}\!\mathit{ef}}
\usepackage{cleveref}
\crefname{originallemma}{original lemma}{original lemmas}
\Crefname{originallemma}{Original lemma}{Original lemmas}
\makeatletter
\newlabel{lemma-consequence-Artin-Rees}{{0EGU}{source}{}{sourcepointer-1}{}}
\newlabel{lemma-consequence-Artin-Rees@cref}{{[originallemma][0][]0EGU}{[1][0][]source}{}{}{}}
\makeatother
\crefname{originallemma}{original lemma}{original lemmas}
\Crefname{originallemma}{Original lemma}{Original lemmas}
\makeatletter
\newlabel{algebra-lemma-completion-complete}{{031C}{source}{}{sourcepointer-2}{}}
\newlabel{algebra-lemma-completion-complete@cref}{{[originallemma][0][]031C}{[1][0][]source}{}{}{}}
\makeatother
\crefname{originallemma}{original lemma}{original lemmas}
\Crefname{originallemma}{Original lemma}{Original lemmas}
\makeatletter
\newlabel{algebra-lemma-NAK}{{00DV}{source}{}{sourcepointer-3}{}}
\newlabel{algebra-lemma-NAK@cref}{{[originallemma][0][]00DV}{[1][0][]source}{}{}{}}
\makeatother
\crefname{originallemma}{original lemma}{original lemmas}
\Crefname{originallemma}{Original lemma}{Original lemmas}
\makeatletter
\newlabel{algebra-lemma-fun}{{05G8}{source}{}{sourcepointer-4}{}}
\newlabel{algebra-lemma-fun@cref}{{[originallemma][0][]05G8}{[1][0][]source}{}{}{}}
\makeatother
\title{Mittag-Leffler property and bounded kernels for the systems of homomorphisms and isomorphisms of the quotients M over I\textasciicircum{}n M}
\author{The Stacks Project Authors}
\date{stacks-project:more-algebra.tex:local-snapshot-20261012; GNU FDL 1.2 or later}
\begin{document}
\maketitle
\setcounter{subsection}{1}
\noindent\textit{Editorial note.} All delivered mathematical text is the author's own, copied from the pinned Stacks Project chapter file more-algebra.tex; no statement or proof step was rewritten, reordered or removed. The counted statement is the labeled lemma environment at source lines 29059-29078 (canonical tag 09BB) and the counted proof is the proof environment at source lines 29080-29121, which immediately follows it and is its own proof.\par\medskip
\noindent\textit{Editorial note.} The printed text at source line 29101 contains the slip by [art part (6) of the lemma where by part (6) of the lemma is meant. It is copied exactly as printed and is not corrected, since the lane delivers the author's own text.\par\medskip
\noindent\textit{Editorial note.} The proof refers four times: once to the Artin-Rees consequence of this chapter and three times to labelled results of another chapter (completeness of completions, Nakayama, and the surjectivity criterion). All are navigation pointers with their roles recorded in the external reference reasons.\par\medskip
\noindent\textit{Editorial note.} The chapter file pinned-original-source.tex carries only an input of the shared preamble on its first line, so no preamble text lives inside the chapter. The retained file sources/stacks-more-algebra/retained-preamble.tex is a mechanical, line-for-line copy of the project's single shared preamble as staged in the snapshot (Xy-pic setup plus the native theorem declarations and mathematical macros). Nothing was added, removed or redefined, and the excerpt serves only as preamble, never as a span coordinate.\par\medskip
\noindent\textit{Editorial note.} The source URL is the source-level Stacks Project URL, as required when more than one card is produced from one source; the per-claim tag is carried by claim\_key and by this note. The counted statement and proof cite no bibliography entry, so this card needs no bibliography.\par\medskip
\noindent\textit{Allowed standard prerequisites (editorial).}\begin{itemize}
\item The Artin-Rees consequence supplies a constant c such that the kernels and cokernels of the comparison maps between Hom(M, N) modulo I\textasciicircum{}n and Hom(M\_n, N\_n) are killed by I\textasciicircum{}c. invoked at proof lines 29094-29098 and used at proof lines 28999-29101. Produces both the constant of part (2) and the Mittag-Leffler property of part (1).
\item For finite modules over a Noetherian ring the completions satisfy M over I\textasciicircum{}n equal to M-completion over I\textasciicircum{}n times the completion, and Hom between the completions is the inverse limit of the Hom groups of the quotients. invoked at proof lines 29102-29109. Gives the identification of the limit of the Hom systems with the group of homomorphisms between the completions, which is part (3).
\item Nakayama's lemma and the resulting surjectivity criterion for endomorphisms of finite modules. invoked at proof lines 29117-29120. Shows that the maps alpha and beta in the final argument are isomorphisms, which yields part (4) for the systems of isomorphisms.
\end{itemize}\medskip
\ifdefined\hypertarget\hypertarget{sourcepointer-1}{}\fi
\noindent\textit{Source reference.} Original lemma 0EGU: sources/stacks-more-algebra/pinned-original-source.tex, source lines 28867--28867. The reference supplies the constant c and the uniform vanishing used for parts (1) and (2), invoked at proof lines 29094-29098; the construction of the complex and the deduction of the four assertions are inside the counted span.\par\medskip
\ifdefined\hypertarget\hypertarget{sourcepointer-2}{}\fi
\noindent\textit{Source reference.} Original lemma 031C: sources/stacks-chow/tags.txt, source lines 3979--3979. The reference records that the completion of a finite module over a Noetherian ring is complete, used at proof lines 29107-29109 to identify the inverse limit of the Hom systems with Hom between the completed modules; the identification argument itself is in the counted span.\par\medskip
\ifdefined\hypertarget\hypertarget{sourcepointer-3}{}\fi
\noindent\textit{Source reference.} Original lemma 00DV: sources/stacks-chow/tags.txt, source lines 512--512. Nakayama's lemma is invoked at proof line 29118 to make the composite of two maps surjective in the argument for the systems of isomorphisms; the surrounding argument that turns surjectivity into bijectivity is inside the counted span.\par\medskip
\ifdefined\hypertarget\hypertarget{sourcepointer-4}{}\fi
\noindent\textit{Source reference.} Original lemma 05G8: sources/stacks-chow/tags.txt, source lines 7052--7052. The reference is the surjectivity criterion for endomorphisms of finite modules, invoked at proof lines 29119-29120 to promote the surjection obtained by Nakayama to an isomorphism; it is a standard pointer and the rest of the isomorphism argument is inside the counted span.\par\medskip
\begin{lemma}
\label{lemma-hom-systems-ML}
Let $A$ be a Noetherian ring. Let $I \subset A$ be an ideal. Let $M$, $N$ be
finite $A$-modules. Set $M_n = M/I^nM$ and $N_n = N/I^nN$. Then
\begin{enumerate}
\item the systems $(\Hom_A(M_n, N_n))$ and $(\text{Isom}_A(M_n, N_n))$
are Mittag-Leffler,
\item there exists a $c \geq 0$ such that the kernels and cokernels of
$$
\Hom_A(M, N)/I^n\Hom_A(M, N) \to \Hom_A(M_n, N_n)
$$
are killed by $I^c$ for all $n$,
\item we have
$\lim \Hom_A(M_n, N_n) =\Hom_A(M, N)^\wedge =
\Hom_{A^\wedge}(M^\wedge, N^\wedge)$
\item $\lim \text{Isom}_A(M_n, N_n) =
\text{Isom}_{A^\wedge}(M^\wedge, N^\wedge)$.
\end{enumerate}
Here ${}^\wedge$ denotes usual $I$-adic completion.
\end{lemma}

\begin{proof}
Note that $\Hom_A(M_n, N_n) = \Hom_A(M, N_n)$. Choose a presentation
$$
A^{\oplus t} \to A^{\oplus s} \to M \to 0
$$
Applying the right exact functor $\Hom_A(-, N)$ we obtain a complex
$$
0 \xrightarrow{\alpha} N^{\oplus s} \xrightarrow{\beta} N^{\oplus t}
$$
whose cohomology in the middle is $\Hom_A(M, N)$ and such that for
$n \geq 0$ the cohomology of
$$
0 \xrightarrow{\alpha_n} N_n^{\oplus s} \xrightarrow{\beta_n} N_n^{\oplus t}
$$
is $\Hom_A(M_n, N_n)$. Let $c \geq 0$ be as in
Lemma \ref{lemma-consequence-Artin-Rees}
for this $A$, $I$, $\alpha$, and $\beta$.
By part (3) of the lemma
we deduce the Mittag-Leffler property for $(\Hom_A(M_n, N_n))$.
The kernel and cokernel of the maps
$\Hom_A(M, N)/I^n\Hom_A(M, N) \to \Hom_A(M_n, N_n)$
are killed by $I^c$ by [art part (6) of the lemma.
We find that $\lim \Hom_A(M_n, N_n) = \Hom_A(M, N)^\wedge$ by
part (2) of the lemma. The equality
$$
\Hom_{A^\wedge}(M^\wedge, N^\wedge) = \lim \Hom_A(M_n, N_n)
$$
follows formally from the fact that $M^\wedge = \lim M_n$ and
$M_n = M^\wedge/I^nM^\wedge$ and the corresponding facts for $N$, see
Algebra, Lemma \ref{algebra-lemma-completion-complete}.

\medskip\noindent
The result for isomorphisms follows from the case of homomorphisms
applied to both $(\Hom(M_n, N_n))$ and $(\Hom(N_n, M_n))$
and the following fact: for $n > m > 0$, if we have maps
$\alpha : M_n \to N_n$ and $\beta : N_n \to M_n$ which
induce an isomorphisms $M_m \to N_m$ and $N_m \to M_m$, then
$\alpha$ and $\beta$ are isomorphisms. Namely, then $\alpha \circ \beta$
is surjective by Nakayama's lemma (Algebra, Lemma \ref{algebra-lemma-NAK})
hence $\alpha \circ \beta$ is an isomorphism by
Algebra, Lemma \ref{algebra-lemma-fun}.
\end{proof}

\end{document}
